\section{Propuesta de Solución}

La propuesta consiste en desarrollar una aplicación web llamada Doble Agenda para el alumnado con discapacidad intelectual del centro Fundación Purísima Concepción de Granada. Su objetivo es ayudar a sus alumnos a organizar sus actividades diarias y escolares, consultar qué tareas tienen pendientes y seguir los pasos necesarios para poder completarlas. Además, permitirá que los profesores organicen las actividades y que las familias puedan también consultar la información del alumno.

\subsection{Funcionalidades que sí se van a desarrollar}
En primer lugar, la aplicación contará con diferentes perfiles de usuario: administrador, profesor, familiar y alumno. Cada uno tendrá acceso a las funciones que le correspondan. Se incluirá la gestión de cuentas y permisos, además de la posibilidad de que los profesores puedan consultar a sus alumnos asignados y que las familias puedan acceder a la información de sus propios hijos.

También, se desarrollará una agenda que permita consultar las actividades del día y de la semana, ver sus detalles e instrucciones y crear actividades personales o escolares. Los profesores podrán asignarlas tanto a alumnos individuales como a grupos de clase. Por su parte, los alumnos podrán indicar cuándo empiezan y terminan una actividad y marcar los pasos que hayan completado cuando una tarea esté dividida en varias partes.

Por último, se incluirán opciones de personalización de la interfaz, como la configuración de colores y la posibilidad de cambiar la foto de perfil. Estas funciones buscan que la aplicación resulte más cómoda y fácil de utilizar, adaptándose así en todo lo posible a las preferencias del alumnado.

Se han priorizado este tipo de funcionalidades porque constituyen la base fundamental de la aplicación, además de que permiten facilitar la organización de las actividades y buscan fomentar la autonomía de los alumnos.

\subsection{Funcionalidades opcionales}
Algunas funcionalidades se han considerado opcionales, ya que, aunque pueden llegar a mejorar la experiencia de uso, no son imprescindibles para que la aplicación cumpla con su objetivo principal. Entre ellas se encuentran la edición o cancelación de actividades, los recordatorios y avisos, el indicador de la hora o reloj, las notificaciones de cancelación y algunas opciones adicionales de personalización, como pueden ser modificar el tamaño del texto o introducir texto mediante la voz.

La incorporación de estas funciones dependerá del tiempo disponible y de cómo avance el desarrollo del proyecto. Al tratarse de un trabajo de la universidad, es necesario compaginarlo con otras asignaturas, prácticas y exámenes, por lo que se ha decidido priorizar las funcionalidades esenciales y básicas, y dejar estas mejoras para el caso de que tengamos tiempo suficiente.

\subsection{Funcionalidades que quedan fuera del proyecto}
En principio, se ha decidido no incluir un canal de comunicación entre usuarios, la sugerencia automática de actividades alternativas cuando se cancela una, ni la posibilidad de que los alumnos valoren las actividades generadas por la aplicación.

El motivo principal por el que hemos decidido no añadir estas funcionalidades es porque ampliarían la complejidad del proyecto y no son realmente necesarias para cubrir las necesidades básicas de organización y seguimiento de actividades. Además, requerirían dedicar más tiempo a su diseño e implementación y, teniendo en cuenta el tiempo disponible del grupo y la carga de trabajo de otras asignaturas con prácticas y exámenes, se ha preferido dedicar el tiempo de desarrollo a las funcionalidades que aportan más valor al objetivo principal de la aplicación Doble Agenda. No obstante, en un futuro, podrían estudiarse como posibles mejoras de la aplicación en otras versiones posteriores.

\subsection{Arquitectura de la solución}
Doble Agenda será una aplicación web accesible desde un navegador, sin necesidad de tener que instalar ningún programa. Se utilizará una arquitectura cliente-servidor, en la que el cliente corresponde a la interfaz con la que interactúan los usuarios, mientras que el servidor se encargará de procesar las peticiones y gestionar la información relacionada con los usuarios, las actividades y las agendas.

La aplicación se basará en una interfaz sencilla, intuitiva y adaptada a las necesidades de los distintos perfiles de usuario, facilitando así la navegación y el acceso a las funcionalidades principales.